\chapter{Plancherel and L2 Fourier Theory}
\label{ch:iii13}

Plancherel theory extends the Fourier transform from integrable functions to a unitary operator on $L^2$. Orthogonality and energy conservation become the main language, allowing transform methods beyond pointwise integral formulas.

\section*{Learning goals}
After this chapter, the reader should be able to:
\begin{itemize}
\item apply Plancherel's theorem to convert \(L^2\) norm or inner-product problems into frequency-space calculations;
\item extend the Fourier transform from an integrable dense subclass to \(L^2\);
\item use Parseval identities to compute energies and orthogonality relations;
\item distinguish \(L^1\)-based pointwise transform statements from \(L^2\)-equivalence-class statements;
\item prove uniqueness results from vanishing Fourier transform in \(L^2\);
\item solve elementary convolution or multiplier problems by exploiting the unitary \(L^2\) transform structure;
\end{itemize}
\section*{Conceptual roadmap}
\[
\boxed{\text{structure}\;\longrightarrow\;\text{estimate}\;\longrightarrow\;\text{limit or representation}\;\longrightarrow\;\text{application}.}
\]

\paragraph{Mathematical setting.}
Use the whole-space convention of III/11 and the complex Hilbert pairing $\langle f,g\rangle=\int f\overline g$. For $f\in L^1\cap L^2$, the completed transform agrees a.e. with the integral transform (approximate simultaneously in both norms). For $f\in L^2$ and $g\in L^1$, Young gives $f*g\in L^2$, and $\widehat{f*g}=\widehat f\widehat g$ in $L^2$, by density from the integrable core.

\section{Parseval identity}
For suitable functions, the inner product is preserved by the Fourier transform.

\section{Plancherel extension}
% VOL03-DEPENDENCY-PREVIEW-III13-SCHWARTZ-DENSE-CORE
\begin{remark}[Dense-core convention]
This chapter uses Schwartz functions as a convenient dense core for the
$L^2$ Fourier transform.  Here $\mathcal S(\mathbb R^n)$ has the local meaning
given in Chapter~11: smooth functions whose polynomially weighted derivatives
are bounded.  Its density in $L^2$ is the standard cutoff-and-mollification
approximation lemma used as an input to the completion argument.  Chapter~14
develops the full Schwartz topology; that later development is not a hidden
prerequisite for the present Plancherel extension.
\end{remark}
Density of Schwartz functions extends the transform uniquely to all of $L^2$.


% BEGIN VOL03-EXPANSION III13-example-01
\begin{example}[Extension by completion]\label{ex:iii13-expand-01}
Choose Schwartz \(f_n\to f\) in \(L^2\). On Schwartz space,
\[
\|\widehat f_n-\widehat f_m\|_2=\|f_n-f_m\|_2,
\]
so the transformed sequence is Cauchy. Define \(\widehat f\) as its \(L^2\)
limit. The isometry also proves independence of the approximating sequence,
giving a unique isometric extension. To obtain surjectivity, also extend the inverse Schwartz transform by the same argument; the two compositions are the identity on the dense core and hence on all of L2.
\end{example}
% END VOL03-EXPANSION III13-example-01
\section{Unitarity}
The $L^2$ transform preserves norm and inner product.

\section{Inverse transform in $L^2$}
Inversion holds in the Hilbert-space sense even when neither function nor transform is integrable.

\section{Convolution in $L^2$ settings}
Combining $L^1$ and $L^2$ hypotheses yields useful convolution bounds and multiplier formulas.

\section{Multipliers}
Bounded frequency multipliers define bounded operators on $L^2$.


% BEGIN VOL03-EXPANSION III13-example-02
\begin{example}[Bounded multiplier]\label{ex:iii13-expand-02}
For
\[
Tf=\mathcal F^{-1}(m\widehat f),\qquad m\in L^\infty,
\]
Plancherel gives
\[
\|Tf\|_2=\|m\widehat f\|_2
\le\|m\|_\infty\|f\|_2.
\]
Thus a frequency-domain sup bound is exactly the operator control needed in
\(L^2\).
\end{example}
% END VOL03-EXPANSION III13-example-02
\section{Energy methods}
Differentiation identities translate Sobolev-type derivative energy into weighted frequency energy.

\section{Approximation}
Frequency cutoffs approximate $L^2$ functions and provide smooth or band-limited regularizations.

% VOL03-NUMERICAL-PLANCHEREL-III13
\begin{remark}[Energy checks for computed transforms]
Plancherel's theorem is an exact unitary identity in \(L^2\).  A computed
transform should be checked against the normalization and discretization
actually used.  Two useful diagnostics are a reconstruction residual and an
energy defect.  In schematic form,
\[
\rho_{\mathcal F}
=
\frac{\|f-\mathcal F^{-1}_{\rm comp}\mathcal F_{\rm comp}f\|_2}{\|f\|_2},
\qquad
\eta_{\rm energy}
=
\frac{\bigl|\|f\|_2^2-\|\widehat f_{\rm comp}\|_2^2\bigr|}{\|f\|_2^2},
\]
where the discrete quadrature/normalization behind both norms must be stated.

A small energy defect checks approximate norm preservation.  A small
reconstruction residual checks approximate inversion.  Neither one alone
certifies adequate spatial/frequency resolution, absence of aliasing, or
pointwise accuracy.  Exact Plancherel identities therefore remain the
mathematical reference, while numerical verification must name the grid,
normalization, truncation, and tolerance used.
\end{remark}


% BEGIN VOL03-EXPANSION III13-example-03
\begin{example}[Frequency cutoffs approximate energy]\label{ex:iii13-expand-03}
Let
\[
\widehat f_R=1_{\{|\xi|\le R\}}\widehat f.
\]
Then
\[
\|f-f_R\|_2^2
=\int_{|\xi|>R}|\widehat f(\xi)|^2\,d\xi\to0.
\]
Band-limited functions therefore approximate arbitrary \(L^2\) functions in
energy, even when pointwise transform formulas are unavailable.
\end{example}
% END VOL03-EXPANSION III13-example-03
\section{Core structural results}

\begin{theorem}[Plancherel theorem]\label{thm:iii13-01}
The Fourier transform extends uniquely to a unitary operator on $L^2(\mathbb R^n)$.
\begin{proof}
For $f,g\in\mathcal S$, inversion and Fubini give $\int\widehat f\,\overline{\widehat g}=\int f\overline g$: absolute integrability follows from $f,\widehat g\in L^1$. Thus the transform is isometric on the dense Schwartz core and extends by completion. The inverse transform has the same isometry property and extends too. Both compositions equal the identity on the core, hence everywhere by continuity. This proves surjectivity as well as norm preservation.
\end{proof}
\end{theorem}

\begin{theorem}[Parseval identity]\label{thm:iii13-02}
For $f,g\in L^2$, $\langle\widehat f,\widehat g\rangle=\langle f,g\rangle$.
\begin{proof}
Use polarization of the norm identity obtained from Plancherel.
\end{proof}
\end{theorem}

\begin{theorem}[$L^2$ multiplier bound]\label{thm:iii13-03}
If $m\in L^\infty$ and $Tf=\mathcal F^{-1}(m\widehat f)$, then $\|Tf\|_2\le\|m\|_\infty\|f\|_2$.
\begin{proof}
Apply Plancherel and the pointwise bound $|m\widehat f|\le\|m\|_\infty|\widehat f|$.
\end{proof}
\end{theorem}

\section{Worked examples}

\begin{example}[Plancherel on a Gaussian]\label{ex:iii13-worked-01}
For \(g(x)=e^{-\pi x^2}\), self-duality gives
\[
\|\widehat g\|_2=\|g\|_2
\]
immediately. Directly, both squared norms are
\(\int_{\mathbb R}e^{-2\pi x^2}dx\), so the identity is visible without any
density argument.
\end{example}

\begin{example}[Extension from a dense class]\label{ex:iii13-worked-02}
Suppose the transform is an isometry on \(\mathcal S\), and
\(f_k\in\mathcal S\) is Cauchy in \(L^2\). Then
\[
\|\widehat f_k-\widehat f_j\|_2=\|f_k-f_j\|_2,
\]
so the transforms are Cauchy as well. Completeness defines the \(L^2\) Fourier
transform as the unique continuous extension.
\end{example}

\begin{example}[Energy of a derivative]\label{ex:iii13-worked-03}
For sufficiently regular \(f\),
\[
\widehat{\partial_jf}=2\pi i\xi_j\widehat f.
\]
Plancherel therefore yields
\[
\|\partial_jf\|_2^2
=4\pi^2\int \xi_j^2|\widehat f(\xi)|^2\,d\xi.
\]
Frequency weighting measures derivative energy.
\end{example}

\begin{example}[Orthogonality turns convolution into an \(L^2\) tool]\label{ex:iii13-worked-04}
If \(g\in L^1\) and \(f\in L^2\), then
\[
\widehat{f*g}=\widehat f\,\widehat g,\qquad
\|\widehat g\|_\infty\le\|g\|_1.
\]
Plancherel gives
\[
\|f*g\|_2\le\|g\|_1\|f\|_2.
\]
\end{example}

\section{Solved dossiers}
Each dossier combines several ideas from the chapter in a longer argument. The aim is to make hypotheses, calculations, and proof strategy explicit.

\begin{problem}[Plancherel on a Gaussian]\label{prob:iii13-dossier-01}
Verify Plancherel's identity for the normalized Gaussian.
\end{problem}
\begin{solution}
Because \(\widehat g=g\), the two \(L^2\) norms are identical; direct
integration gives the same quantity on each side.
\end{solution}

\begin{problem}[Extension from a dense class]\label{prob:iii13-dossier-02}
Explain how an isometry on a dense subspace extends uniquely to all of \(L^2\).
\end{problem}
\begin{solution}
Approximate \(f\) by a dense-class sequence. Isometry makes the transformed
sequence Cauchy, and the limit is independent of the approximating sequence.
\end{solution}

\begin{problem}[Energy of a derivative]\label{prob:iii13-dossier-03}
Use Plancherel to express \(\|\partial_jf\|_2\) in frequency variables.
\end{problem}
\begin{solution}
Insert the derivative multiplier \(2\pi i\xi_j\) and use equality of
\(L^2\) norms.
\end{solution}

\begin{problem}[Orthogonality turns convolution into an \(L^2\) tool]\label{prob:iii13-dossier-04}
Derive the \(L^1*L^2\to L^2\) convolution estimate using Fourier analysis.
\end{problem}
\begin{solution}
Use the convolution theorem, the \(L^1\to L^\infty\) transform bound, and
Plancherel.
\end{solution}

\begin{problem}[Parseval polarization]\label{prob:iii13-dossier-05}
How can inner-product preservation be recovered from norm preservation?
\end{problem}
\begin{solution}
Use the polarization identity; an isometry on a complex or real Hilbert
space preserves the associated inner product.
\end{solution}

\begin{problem}[Dense-class independence]\label{prob:iii13-dossier-06}
Why must the \(L^2\) transform limit be independent of the approximating Schwartz sequence?
\end{problem}
\begin{solution}
Two approximating sequences differ by an \(L^2\)-null sequence; isometry
makes their transforms differ by an \(L^2\)-null sequence too.
\end{solution}

\begin{problem}[Inverse transform in \(L^2\)]\label{prob:iii13-dossier-07}
What does Plancherel imply about invertibility of the Fourier transform on \(L^2\)?
\end{problem}
\begin{solution}
The unitary extension has inverse given by the inverse Fourier transform,
defined in the same \(L^2\) sense.
\end{solution}

\begin{problem}[Frequency truncation]\label{prob:iii13-dossier-08}
Why do low-pass truncations converge in \(L^2\)?
\end{problem}
\begin{solution}
Multiplying \(\widehat f\) by indicators of expanding frequency balls
converges to \(\widehat f\) in \(L^2\); Plancherel transfers this convergence
back to physical space.
\end{solution}

\begin{problem}[Plancherel theorem proof dossier]\label{prob:iii13-dossier-09}
Prove the following structural statement and identify the step where its main hypothesis is used: The Fourier transform extends uniquely to a unitary operator on $L^2(\mathbb R^n)$.
\end{problem}
\begin{solution}
For $f,g\in\mathcal S$, inversion and Fubini give $\int\widehat f\,\overline{\widehat g}=\int f\overline g$: absolute integrability follows from $f,\widehat g\in L^1$. Thus the transform is isometric on the dense Schwartz core and extends by completion. The inverse transform has the same isometry property and extends too. Both compositions equal the identity on the core, hence everywhere by continuity. This proves surjectivity as well as norm preservation. This proof also explains why the stated hypothesis is structural rather than cosmetic.
\end{solution}

\begin{problem}[Parseval identity proof dossier]\label{prob:iii13-dossier-10}
Prove the following structural statement and identify the step where its main hypothesis is used: For $f,g\in L^2$, $\langle\widehat f,\widehat g\rangle=\langle f,g\rangle$.
\end{problem}
\begin{solution}
Use polarization of the norm identity obtained from Plancherel. This proof also explains why the stated hypothesis is structural rather than cosmetic.
\end{solution}

\begin{problem}[$L^2$ multiplier bound proof dossier]\label{prob:iii13-dossier-11}
Prove the following structural statement and identify the step where its main hypothesis is used: If $m\in L^\infty$ and $Tf=\mathcal F^{-1}(m\widehat f)$, then $\|Tf\|_2\le\|m\|_\infty\|f\|_2$.
\end{problem}
\begin{solution}
Apply Plancherel and the pointwise bound $|m\widehat f|\le\|m\|_\infty|\widehat f|$. This proof also explains why the stated hypothesis is structural rather than cosmetic.
\end{solution}

\begin{problem}[Chapter synthesis]\label{prob:iii13-dossier-12}
Connect the main definitions and structural theorems of this chapter into one proof strategy. Explain what object is controlled, what estimate or closure principle is used, and what conclusion becomes available.
\end{problem}
\begin{solution}
Plancherel theory extends the Fourier transform from integrable functions to a unitary operator on $L^2$. Orthogonality and energy conservation become the main language, allowing transform methods beyond pointwise integral formulas. The recurring strategy is to encode the object in the chapter's natural measurable, normed, Fourier, or distributional structure, establish the controlling estimate, and then pass to the desired limit or representation.
\end{solution}
\section{Exercises with complete solutions}
The shorter exercise layer remains separate from the solved dossiers.

\begin{exercise}\label{exr:iii13-01}
State and apply the central principle behind Parseval identity. Give one immediate consequence in the setting of this chapter.
\end{exercise}
\begin{hint}
Use Plancherel to move the inner product to frequency space before estimating the transformed expression.
\end{hint}
\begin{solution}
For suitable functions, the inner product is preserved by the Fourier transform. As an immediate consequence, the same principle can be inserted into later convergence, approximation, transform, or weak-form arguments whenever its hypotheses are verified.
\end{solution}

\begin{exercise}\label{exr:iii13-02}
State and apply the central principle behind Plancherel extension. Give one immediate consequence in the setting of this chapter.
\end{exercise}
\begin{hint}
Define the \(L^2\) transform by density from Schwartz or \(L^1\cap L^2\) functions, not by assuming the integral exists pointwise.
\end{hint}
\begin{solution}
Density of Schwartz functions extends the transform uniquely to all of $L^2$. As an immediate consequence, the same principle can be inserted into later convergence, approximation, transform, or weak-form arguments whenever its hypotheses are verified.
\end{solution}

\begin{exercise}\label{exr:iii13-03}
State and apply the central principle behind Unitarity. Give one immediate consequence in the setting of this chapter.
\end{exercise}
\begin{hint}
Parseval is the inner-product form of Plancherel; polarize the norm identity if necessary.
\end{hint}
\begin{solution}
The $L^2$ transform preserves norm and inner product. As an immediate consequence, the same principle can be inserted into later convergence, approximation, transform, or weak-form arguments whenever its hypotheses are verified.
\end{solution}

\begin{exercise}\label{exr:iii13-04}
State and apply the central principle behind Inverse transform in $L^2$. Give one immediate consequence in the setting of this chapter.
\end{exercise}
\begin{hint}
Two \(L^2\) functions that differ on a null set represent the same transform class, so avoid pointwise claims unless extra regularity is present.
\end{hint}
\begin{solution}
Inversion holds in the Hilbert-space sense even when neither function nor transform is integrable. As an immediate consequence, the same principle can be inserted into later convergence, approximation, transform, or weak-form arguments whenever its hypotheses are verified.
\end{solution}

\begin{exercise}\label{exr:iii13-05}
State and apply the central principle behind Convolution in $L^2$ settings. Give one immediate consequence in the setting of this chapter.
\end{exercise}
\begin{hint}
If the Fourier transform vanishes in \(L^2\), Plancherel forces the original function to have zero \(L^2\) norm.
\end{hint}
\begin{solution}
Combining $L^1$ and $L^2$ hypotheses yields useful convolution bounds and multiplier formulas. As an immediate consequence, the same principle can be inserted into later convergence, approximation, transform, or weak-form arguments whenever its hypotheses are verified.
\end{solution}

\begin{exercise}\label{exr:iii13-06}
State and apply the central principle behind Multipliers. Give one immediate consequence in the setting of this chapter.
\end{exercise}
\begin{hint}
Multipliers bounded in frequency define bounded \(L^2\) operators; estimate by the essential supremum of the multiplier.
\end{hint}
\begin{solution}
Bounded frequency multipliers define bounded operators on $L^2$. As an immediate consequence, the same principle can be inserted into later convergence, approximation, transform, or weak-form arguments whenever its hypotheses are verified.
\end{solution}

\begin{exercise}\label{exr:iii13-07}
State and apply the central principle behind Energy methods. Give one immediate consequence in the setting of this chapter.
\end{exercise}
\begin{hint}
For convolution in \(L^2\), verify the other factor provides enough integrability to make the product formula legitimate.
\end{hint}
\begin{solution}
Differentiation identities translate Sobolev-type derivative energy into weighted frequency energy. As an immediate consequence, the same principle can be inserted into later convergence, approximation, transform, or weak-form arguments whenever its hypotheses are verified.
\end{solution}

\begin{exercise}\label{exr:iii13-08}
State and apply the central principle behind Approximation. Give one immediate consequence in the setting of this chapter.
\end{exercise}
\begin{hint}
Check normalization constants by applying Plancherel to a Gaussian, where both sides can be computed independently.
\end{hint}
\begin{solution}
Frequency cutoffs approximate $L^2$ functions and provide smooth or band-limited regularizations. As an immediate consequence, the same principle can be inserted into later convergence, approximation, transform, or weak-form arguments whenever its hypotheses are verified.
\end{solution}


% BEGIN VOL03-EXPANSION III13-exercises-01
\section{Graded supplementary exercises}
These exercises extend the preserved foundational set and are graded by mathematical role.

\subsection*{Standard computations}

\begin{exercise}[Energy]\label{exr:iii13-expand-01}
If ||f||\_2=5, what is ||hat f||\_2?
\end{exercise}
\begin{hint}
Use Plancherel.
\end{hint}
\begin{solution}
5.
\end{solution}

\begin{exercise}[Multiplier]\label{exr:iii13-expand-02}
If ||m||\_infinity=3, bound ||Tf||\_2.
\end{exercise}
\begin{hint}
Use multiplier bound.
\end{hint}
\begin{solution}
At most 3||f||\_2.
\end{solution}

\begin{exercise}[Parseval]\label{exr:iii13-expand-03}
What happens to inner products under Fourier transform in L2?
\end{exercise}
\begin{hint}
Use Parseval.
\end{hint}
\begin{solution}
They are preserved.
\end{solution}

\begin{exercise}[Cutoff error]\label{exr:iii13-expand-04}
Express the L2 error of a sharp frequency cutoff.
\end{exercise}
\begin{hint}
Use Plancherel.
\end{hint}
\begin{solution}
It is the L2 mass of hat f outside the cutoff.
\end{solution}

\begin{exercise}[Zero transform]\label{exr:iii13-expand-05}
If hat f=0 in L2, what is f?
\end{exercise}
\begin{hint}
Use unitarity.
\end{hint}
\begin{solution}
f=0 almost everywhere.
\end{solution}

\subsection*{Proofs}

\begin{exercise}[Extension uniqueness]\label{exr:iii13-expand-06}
Why is the L2 extension unique?
\end{exercise}
\begin{hint}
Schwartz space is dense and transform is continuous.
\end{hint}
\begin{solution}
Two continuous extensions agreeing on a dense set agree everywhere.
\end{solution}

\begin{exercise}[Parseval from norm]\label{exr:iii13-expand-07}
Derive inner-product preservation from norm preservation.
\end{exercise}
\begin{hint}
Use polarization.
\end{hint}
\begin{solution}
Polarization reconstructs the inner product from norms.
\end{solution}

\begin{exercise}[Multiplier proof]\label{exr:iii13-expand-08}
Prove the multiplier bound.
\end{exercise}
\begin{hint}
Use Plancherel and pointwise multiplier bound.
\end{hint}
\begin{solution}
Norm Tf equals norm m hat f, bounded by ||m||infinity ||f||2.
\end{solution}

\begin{exercise}[Cutoff convergence]\label{exr:iii13-expand-09}
Prove sharp cutoffs converge in L2.
\end{exercise}
\begin{hint}
Use integrable tail of |hat f|\textasciicircum{}2.
\end{hint}
\begin{solution}
The tail integral tends to zero.
\end{solution}

\subsection*{Counterexamples and hypothesis testing}

\begin{exercise}[L2 not L1]\label{exr:iii13-expand-10}
Why may an L2 transform lack a pointwise integral formula?
\end{exercise}
\begin{hint}
L2 does not imply L1 on the whole real line.
\end{hint}
\begin{solution}
The transform is then defined as an equivalence class by completion.
\end{solution}

\begin{exercise}[Unbounded multiplier]\label{exr:iii13-expand-11}
Why can an unbounded multiplier fail to be a bounded L2 operator?
\end{exercise}
\begin{hint}
Concentrate spectral mass where m is large.
\end{hint}
\begin{solution}
Operator norms would become arbitrarily large.
\end{solution}

\begin{exercise}[Pointwise identity]\label{exr:iii13-expand-12}
Why are L2 transform identities generally almost-everywhere statements?
\end{exercise}
\begin{hint}
L2 elements are equivalence classes.
\end{hint}
\begin{solution}
Point values are not canonical.
\end{solution}

\subsection*{Applications and investigations}

\begin{exercise}[Energy filter]\label{exr:iii13-expand-13}
How does a spectral mask change energy?
\end{exercise}
\begin{hint}
Use Plancherel.
\end{hint}
\begin{solution}
Output energy is the L2 energy of the masked spectrum.
\end{solution}

\begin{exercise}[Derivative energy]\label{exr:iii13-expand-14}
Why does gradient energy become weighted spectral energy?
\end{exercise}
\begin{hint}
Fourier derivative multiplies by frequency.
\end{hint}
\begin{solution}
The weight is proportional to |xi|\textasciicircum{}2.
\end{solution}

\subsection*{Challenge problems}

\begin{exercise}[Projection]\label{exr:iii13-expand-15}
Show a sharp frequency cutoff is an orthogonal projection.
\end{exercise}
\begin{hint}
Its multiplier is a real indicator with m\textasciicircum{}2=m.
\end{hint}
\begin{solution}
The operator is self-adjoint and idempotent.
\end{solution}

\begin{exercise}[Sobolev preview]\label{exr:iii13-expand-16}
Express L2 plus gradient energy in frequency terms.
\end{exercise}
\begin{hint}
Use Plancherel and derivative rule.
\end{hint}
\begin{solution}
It is an integral of (1+4 pi\textasciicircum{}2 |xi|\textasciicircum{}2)|hat f|\textasciicircum{}2.
\end{solution}
% END VOL03-EXPANSION III13-exercises-01
\begin{exercise}[Resolved and unresolved Fourier energy]\label{exr:iii13-numerical-energy-01}
Assume \(\|f\|_2=1\).  Exact Fourier information shows that the modes retained
by a cutoff contain squared \(L^2\) energy \(0.96\).

\begin{enumerate}
\item Compute the exact \(L^2\) truncation error associated with the omitted
frequency tail.
\item A numerical transform/reconstruction instead gives a measured relative
\(L^2\) residual \(0.23\).  Explain why the difference between \(0.23\) and the
exact tail error should not automatically be called truncation error.
\end{enumerate}
\end{exercise}
\begin{hint}
Use orthogonality/Plancherel for the exact cutoff:
\(\|f-f_R\|_2^2=1-0.96\).  Then separate exact tail energy from coefficient,
sampling, quadrature, and reconstruction error.
\end{hint}
\begin{solution}
The omitted energy is \(1-0.96=0.04\), so the exact cutoff error is
\[
\|f-f_R\|_2=\sqrt{0.04}=0.2.
\]
The numerical residual \(0.23\) is larger.  That residual includes whatever
errors are present in the computed coefficients, transform normalization,
sampling, quadrature, inverse transform, and any other discretization step.
Only the exact projection identity isolates the omitted Fourier tail as
\(0.2\).  The additional discrepancy must therefore be diagnosed separately
rather than silently attributed to truncation.
\end{solution}

\section*{Chapter summary}
The chapter now supplies a canonical theorem-and-dossier layer for subsequent measure, Fourier, distributional, Sobolev, and PDE arguments.
